\section{Анализ предметной области}
\subsection{Регистрация и учёт данных пользователей: эволюция подходов}

Исторически учёт «пользователей» и связанных с ними данных начинался с бумажных журналов и картотек: кто обратился, что сделано, кем выполнено, когда. С появлением вычислительной техники в 1970–1980-х годах первые автоматизированные системы фокусировались в основном на учёте фактов — склад, бухгалтерия, справочники; данные о сотрудниках и их ролях часто оставались на уровне учётных записей операционной системы или отдельных учётных модулей. Понятие «пользователь системы» как субъекта с определёнными правами и ролью получило широкое распространение с развитием многопользовательских СУБД и корпоративных информационных систем в 1990-х годах.

В 1990-х годах с распространением персональных компьютеров и локальных сетей стали типичными приложения с явным разграничением доступа: учётные записи, роли (администратор, оператор, исполнитель и т.д.), привязка операций к конкретному пользователю. Данные хранились в локальных базах (dBase, FoxPro, позже Access, SQL Server), доступ — только с рабочих мест внутри предприятия. Регистрация данных пользователей означала уже не только ведение справочника сотрудников, но и фиксацию того, кто какую операцию выполнил, что принципиально важно для ответственности, аудита и разбора спорных ситуаций.

В 2000-х годах с развитием веб-технологий и интернета появились облачные и гибридные решения. Регистрация данных пользователей стала включать: единую аутентификацию, профили с ролями, привязку заявок и операций к пользователям, журналы изменений (аудит). Одновременно усилились требования к безопасности персональных данных и к разграничению доступа на уровне приложения и базы данных — так, чтобы пользователь не мог ни прочитать, ни изменить данные, на которые у него нет прав.

В наши дни типичная программно-информационная система регистрации данных пользователей обеспечивает: ведение учётных записей и ролей; регистрацию и учёт связанных сущностей (клиенты, объекты, заявки, операции); регламентированный жизненный цикл операций с проверкой прав; фиксацию изменений для аудита. Такая система снижает ошибки, ускоряет процессы и даёт руководству актуальную и прозрачную картину. Конкретная отрасль или сценарий (например, сервисный центр, автосервис, учебное заведение) выбирается как <<пример реализации>> — предметная область при этом остаётся общей: регистрация и учёт данных пользователей.

Таким образом, эволюция от бумажного и локального учёта к централизованным, ролевым и безопасным программным системам формирует общую предметную область данной разработки; выбор конкретного примера (в данном случае — CRM учёта на примере автосервиса) позволяет наглядно реализовать все ключевые аспекты ПИС регистрации данных пользователей.

\subsection{Характеристики предметной области и выбор примера реализации}

Предметная область — <<регистрация и учёт данных пользователей>> в широком смысле: пользователи системы (сотрудники с ролями и правами), связанные с ними сущности (клиенты, объекты учёта, заявки) и фиксация изменений (аудит). Чтобы реализовать такую систему в рамках дипломного проекта в полном объёме, было принято решение реализовать <<CRM-систему учёта пользователей (и не только) на примере автосервиса>>. Автосервис как пример даёт наглядные сущности: клиенты, автомобили, заявки на обслуживание, менеджеры и мастера с разными ролями, регламентированный жизненный цикл заявок — при этом архитектура и принципы остаются переносимыми на другие отрасли (сервисные центры, приёмные комиссии, службы заявок и т.д.).

В выбранном примере реализации ключевые сущности выглядят следующим образом.

<<Пользователи системы>> — сотрудники, имеющие учётные записи и роли: администратор (полный доступ, в том числе управление пользователями и назначение ролей), менеджер (заявки, клиенты, автомобили, назначение исполнителей, просмотр журнала изменений), мастер/техник (просмотр заявок, взятие в работу, смена статуса назначенных ему заявок). Разграничение доступа обеспечивается на уровне приложения (маршруты, видимость кнопок и разделов) и на уровне базы данных (политики RLS), так что пользователь не может обойти ограничения ни через интерфейс, ни через прямой доступ к API.

<<Клиенты>>  характеризуются учётными данными: ФИО, телефон, электронная почта. С клиентом связываются один или несколько объектов учёта (в примере автосервиса — автомобили). Данные клиентов используются при создании заявок и в отчётности; регистрация и актуализация этих данных входит в зону ответственности пользователей с соответствующей ролью.

<<Объекты учёта>> (в примере — <<автомобили>>): привязаны к клиентам, описываются набором атрибутов (марка, модель, год, госномер, при необходимости VIN). Один объект может фигурировать в нескольких заявках за период эксплуатации. Ведение справочника объектов и привязка к клиентам — часть регистрации данных в системе.

<<Заявки (операции)>> — центральная сущность, связывающая пользователей, клиентов и объекты. У каждой заявки есть клиент, объект (автомобиль), описание работ, статус жизненного цикла (черновик, отправлена, назначена исполнителю, в работе, приостановлена, завершена, отменена), менеджер, создавший заявку, и при необходимости назначенный исполнитель (мастер). Переходы между статусами регламентированы: например, из черновика заявку можно только отправить или отменить; назначить мастера можно только для отправленной заявки. Бизнес-логика переходов и прав доступа по ролям должна быть явной и единообразной — это обеспечивает предсказуемость системы и удобство сопровождения.

<<Журнал изменений (аудит)>> фиксирует, кто, когда и какие изменения внёс в заявки, клиентов, объекты и пользователей (в том числе смену ролей). Это неотъемлемая часть регистрации данных пользователей: система не только хранит текущее состояние, но и оставляет след действий пользователей для прозрачности и разбора спорных ситуаций.

Гармоничная работа системы достигается за счёт чёткого разделения ответственности по ролям, явного жизненного цикла заявок (workflow) и единого источника правды о данных в БД с защитой на уровне RLS. Выбор автосервиса как примера не сужает тему диплома до «системы для автосервиса» — тема остаётся «ПИС регистрации данных пользователей», а автосервис выступает наглядной предметной областью для реализации.

\subsection{Анализ существующего программного обеспечения}

Существует множество программных решений, так или иначе связанных с регистрацией данных пользователей и учётом заявок, клиентов и операций — от локальных десктоп-приложений до облачных CRM и отраслевых систем. Ниже рассмотрены типичные подходы и аналоги, чтобы обосновать выбор архитектуры и функциональности разрабатываемой системы применительно к теме «ПИС регистрации данных пользователей» и примеру CRM на базе автосервиса.

<<Универсальные CRM (Salesforce, Битрикс24, AmoCRM и аналоги)>> gозволяют вести клиентов, сделки, задачи и пользователей с ролями. Плюсы: широкая настройка, интеграции, масштабируемость. Минусы для учебного и дипломного проекта: нет изначальной отраслевой модели «заявка — клиент — объект — статус — исполнитель» в нужном виде, требуется существенная кастомизация; высокая стоимость и сложность внедрения для небольших организаций; закрытый или сложный код не позволяет наглядно продемонстрировать архитектуру регистрации данных и разграничения доступа. Для диплома важнее прозрачная архитектура, контроль над кодом и данными и возможность показать именно механизмы регистрации и учёта данных пользователей.

<<Специализированные отраслевые системы>> (например, для автосервисов: 1С:Автосервис, АвтоГРАФ и др.). Заточены под заказ-наряды, склад запчастей, учёт работ. Плюсы: богатая отраслевая функциональность. Минусы: как правило, платные, требуют обучения и нередко привязки к конкретной платформе; внутренняя реализация регистрации данных и RLS часто скрыта; для дипломного проекта важнее наглядность архитектуры, явный workflow и возможность доработки без лицензионных ограничений.

<<Веб-приложения на современном стеке (Vue, React) + Backend-as-a-Service (Supabase, Firebase)>>. Современный подход: SPA на фронтенде, база данных, аутентификация и API предоставляются облачным сервисом. Плюсы: быстрый старт, встроенная аутентификация и механизмы RLS, возможность сосредоточиться на бизнес-логике регистрации данных и UX; открытая схема БД и политики RLS позволяют наглядно показать, как реализована защита данных на уровне строк. Минусы: зависимость от провайдера BaaS — для дипломного проекта приемлемо. Такой вариант оптимален: демонстрируется полноценная клиент-серверная архитектура, регистрация и учёт пользователей и связанных сущностей, роли, workflow и безопасность данных на уровне приложения и БД.

На основе этого анализа для данной разработки выбран подход: <<Vue 3 (Composition API) + TypeScript на фронтенде, Supabase (PostgreSQL, Auth, RLS) на бэкенде>>, с модульной доменной структурой и явным workflow заявок. Это обеспечивает соответствие теме «ПИС регистрации данных пользователей», наглядность для защиты, возможность показать механизмы разграничения доступа и аудита, а также последующее расширение системы (новые роли, статусы, отчёты, перенос примера на другую отрасль).

\subsection{Перспективы и направления развития}

Анализ предметной области и аналогов показывает следующие направления развития как в рамках примера автосервиса, так и в общем контексте ПИС регистрации данных пользователей.

<<Расширение ролевой модели>>.  Например, введение роли «клиент» с доступом только к своим заявкам через личный кабинет — это типичный сценарий для многих отраслей (сервисные центры, образовательные порталы, медицинские записи). Регистрация данных пользователей при этом дополняется учётом внешних пользователей (клиентов) с ограниченными правами.

<<Уведомления>>, Письма или push-уведомления при смене статуса заявки повышают информированность клиентов и исполнителей; логика остаётся в рамках регистрации событий и привязки к пользователям.

<<Отчёты и экспорт>>.  Отчёты по заявкам, клиентам, выручке, загрузке исполнителей; экспорт в Excel/PDF — расширяют аналитические возможности без изменения ядра системы регистрации данных.

<<Интеграции>>.  Интеграция с 1С, телефонией, SMS-рассылками позволяет использовать зарегистрированные данные пользователей и заявок во внешних системах при сохранении единого источника правды в разрабатываемой ПИС.

<<Мобильное приложение или PWA>>  для исполнителей (например, мастеров «в яме») — тот же набор сущностей и ролей, другой канал доступа; архитектура с API-слоем и RLS допускает добавление мобильного клиента без дублирования бизнес-логики.

<<Real-time обновления>> через подписки Supabase (или аналогов) для актуального отображения изменений на экранах нескольких пользователей — особенно важно в сценариях с совместной работой над заявками.

Разрабатываемая система закладывает фундамент для таких расширений: модульная доменная структура, единый API-слой, RLS и явный workflow позволяют добавлять новые домены и сценарии без ломки существующей логики регистрации и учёта данных пользователей.

\subsection{Выводы по анализу предметной области}

В ходе анализа предметной области были рассмотрены особенности регистрации и учёта данных пользователей, а также прослежена эволюция подходов к организации подобных систем: от бумажных журналов и локальных баз данных до современных веб-ориентированных решений с централизованным хранением информации и развитой системой разграничения прав доступа. Установлено, что в настоящее время программно-информационные системы регистрации данных пользователей должны обеспечивать не только хранение и обработку сведений, но и идентификацию субъектов, выполняющих операции, контроль полномочий пользователей и фиксацию истории изменений данных.

В качестве примера реализации предметной области была выбрана CRM-система автосервиса, поскольку данный сценарий позволяет наглядно продемонстрировать все ключевые механизмы, характерные для программно-информационных систем регистрации данных пользователей. Использование таких сущностей, как пользователи, клиенты, автомобили и заявки на обслуживание, а также наличие регламентированного жизненного цикла заявок позволяют реализовать полный цикл регистрации, обработки и контроля данных. При этом разработанные подходы и архитектурные решения не ограничиваются рамками автосервисной деятельности и могут быть адаптированы для применения в других предметных областях, связанных с обработкой пользовательских данных и управлением обращениями.

Проведённый анализ существующих программных продуктов показал, что универсальные CRM-системы и специализированные отраслевые решения обладают широкими функциональными возможностями, однако зачастую характеризуются высокой стоимостью внедрения, избыточностью функционала или недостаточной гибкостью при адаптации к конкретным требованиям организации. Использование электронных таблиц и простых локальных решений, напротив, не обеспечивает необходимый уровень безопасности, совместной работы пользователей и контроля изменений данных.

На основании проведённого анализа обоснована целесообразность разработки собственной программно-информационной системы регистрации данных пользователей в виде веб-приложения, реализованного с использованием современных технологий Vue 3, TypeScript и Supabase. Выбранный подход позволяет обеспечить централизованное хранение данных, ролевое разграничение доступа, поддержку регламентированных бизнес-процессов, ведение журнала аудита и возможность дальнейшего расширения функциональности системы. Полученные результаты анализа предметной области послужили основой для формирования требований к разрабатываемой системе и последующего выполнения этапов технического задания и проектирования.


